\section{Data}\label{sec:data}

We use lap-level timing for the 2022 to 2026 seasons from the FastF1 Python package \citep[version 3.8.3]{schaefer2026fastf1}, covering practice sessions (FP1 to FP3), sprint sessions and races.
Each lap carries its time, driver, stint, tyre compound label, tyre age in laps (which for a used set includes earlier running), track status and start time.
A \emph{run} is one driver's continuous spell on one set of tyres.
Car telemetry for the check in Appendix~\ref{app:tele} comes from the OpenF1 interface \citep{openf12026}.
For 2026 we use the 13 races completed when the data were collected (Table~\ref{tab:data}).

\begin{table}[t]
\centering\small
\begin{tabular}{@{}lrrr@{}}
\toprule
 & Events or sessions & Laps & Runs \\
\midrule
\multicolumn{4}{@{}l}{\textit{Race fit by season (long-run race laps)}} \\
2022 & 18 & 13{,}637 & 1{,}038 \\
2023 & 16 & 13{,}536 & 860 \\
2024 & 18 & 15{,}942 & 857 \\
2025 & 17 & 15{,}279 & 872 \\
2026 & 13 & 11{,}039 & 820 \\
\midrule
\multicolumn{4}{@{}l}{\textit{2026 clean laps by session type}} \\
Practice 1 & 14 & 4{,}795 &  \\
Practice 2 & 9 & 3{,}039 &  \\
Practice 3 & 9 & 1{,}828 &  \\
Sprint & 5 & 1{,}894 &  \\
Race & 13 & 12{,}300 &  \\
All sessions & 14 & 23{,}856 &  \\
\midrule
\multicolumn{4}{@{}l}{\textit{2026 laps entering the degradation fits}} \\
Race, long runs & 13 & 11{,}039 & 820 \\
Practice and sprint, race-simulation runs & 14 & 2{,}815 & 300 \\
\bottomrule
\end{tabular}

\caption{Data used. Italic rows are panel headings. Race-fit laps are green-flag laps in runs of at least five laps. Practice and sprint runs are those classified as race simulations (Section~\ref{sec:data}). Season rows use hard, medium and soft labels. The pooled 2022--26 fit of Section~\ref{sec:rates} uses physical compounds C1 to C5 and so omits the C6 soft tyre of the 2025 Monaco race, leaving \PooledNLaps{} laps in \PooledNRuns{} runs.}
\label{tab:data}
\end{table}

\medskip\noindent\textit{Lap and run selection.}\ We keep timed laps that are not in-laps or out-laps, were run entirely under green flags and are marked accurate by FastF1.
A long run is a sequence of at least five consecutive such laps on one set of tyres.
In races we drop laps more than 2.5 standard deviations from the run median, which removes lock-ups and incidents.
In practice we also drop laps slower than 1.10 times the driver's best lap (cool-down and preparation laps), split runs at gaps of more than three laps, discard the first lap of each run as warm-up, and keep only runs that look like race simulations: median lap at least 1.5\,s off the session best and a lap-time standard deviation of at most 2.5\,s.
Sprint sessions count as race running.
A race group needs at least four cars on the lap.

\medskip\noindent\textit{Compounds.}\ Pirelli nominates three of its five dry compounds (C1, the hardest, to C5, the softest) for each weekend and labels them hard, medium and soft within that weekend, so a label denotes different rubber at different circuits.
Harder compounds are meant to degrade more slowly \citep[pp.~2, 9]{cappelloHoegh2026}.
We map labels to physical compounds using nominations compiled from Pirelli's published releases, because no timing feed carries them; the per-event table with sources is in the supplementary data.
C numbers are ordinal within a season and the range was re-specified for 2025, so the pooled circuit analysis of Section~\ref{sec:rates} is repeated without 2025 and on the \SensVerifiedNEvents{} races whose nominations two sources confirm.
The other 2022 to 2025 analyses use the labels.
Because Pirelli nominates by circuit, compounds differ by circuit, which limits what any pooled fit can say about compounds (Section~\ref{sec:power}).

\medskip\noindent\textit{Traffic.}\ We compute the gap to the car ahead at the start of each lap among cars on the same lap number.
A car within 2\,s is in dirty air, and the traffic covariate is the shortfall of the gap from 2\,s, with gaps capped at 10\,s.
Cars lapped by the car ahead are invisible to this measure.
Omitting traffic biases the tyre-age slope towards zero, because fresh tyres often rejoin in traffic.
